% 题证摘录；来源与整理边界见相同编号分题文件。
\begin{definition}[Weierstrass polynomial]
Let $U\subset\C^{n-1}$ be open, $0\in U$, and let $z'\in\C^{n-1}$ denote the coordinates. Suppose $P\in\sO(U)[z_n]$ is a monic polynomial of degree $k\geq0$,
\[
P(z',z_n)=z_n^k+\sum_{\ell=0}^{k-1}c_\ell(z')z_n^\ell,
\]
where $c_\ell$ are holomorphic functions defined on $U$ such that $c_\ell(0)=0$ for all $\ell$. Then $P$ is called a Weierstrass polynomial of degree $k$.
\end{definition}
\begin{theorem}[Weierstrass preparation theorem]
Suppose $f\in\sO(U)$ for an open $U\subset\C^{n-1}\times\C$, where $0\in U$, and $f(0)=0$. Suppose $z_n\mapsto f(0,z_n)$ is not identically zero near the origin and its order of vanishing at the origin is $k\geq1$. Then there exists an open polydisc $V=V'\times D\subset\C^{n-1}\times\C$ with $0\in V\subset U$, a unique $u\in\sO(V)$, $u(z)\neq0$ for all $z\in V$, and a unique Weierstrass polynomial $P$ of degree $k$ with coefficients holomorphic in $V'$ such that
\[
f(z',z_n)=u(z',z_n)P(z',z_n),
\]
and such that all $k$ zeros (counting multiplicity) of $z_n\mapsto P(z',z_n)$ lie in $D$ for all $z'\in V'$.
\end{theorem}
\begin{proof}
There exists a small disc $D\subset\C$ centered at zero such that $\{0\}\times\widebar D\subset U$ and such that $f(0,z_n)\neq0$ for $z_n\in\widebar D\setminus\{0\}$. By continuity of $f$, there is a small polydisc $V=V'\times D$ such that $\widebar V\subset U$ and $f$ is not zero on $V'\times\partial D$. See \figureref{fig:wpreppoly} for the setup. We will consider the zeros of $z_n\mapsto f(z',z_n)$ for $z'\in V'$. See \figureref{fig:wprepdiv}.

By the one-variable argument principle (\thmref{thm:onevarargprinc}), the number of zeros (with multiplicity) of $z_n\mapsto f(z',z_n)$ in $D$ is
\[
\frac1{2\pi i}\int_{\partial D}\frac{\frac{\partial f}{\partial z_n}(z',\zeta)}{f(z',\zeta)}\,d\zeta.
\]
As $f(z',\zeta)$ does not vanish when $z'\in V'$ and $\zeta\in\partial D$, the expression above is a continuous integer-valued function of $z'\in V'$. The expression is equal to $k$ when $z'=0$, and so it is equal to $k$ for all $z'\in V'$.
Write the zeros of $z_n\mapsto f(z',z_n)$ as $\alpha_1(z'),\ldots,\alpha_k(z')$, including multiplicity. The zeros are not ordered in any particular way---pick \emph{some} ordering for every $z'$. Write
\[
P(z',z_n)=\prod_{\ell=1}^k\bigl(z_n-\alpha_\ell(z')\bigr)=z_n^k+c_{k-1}(z')z_n^{k-1}+\cdots+c_0(z').
\]
For a fixed $z'$, $P$ (and thus the coefficients $c_0,\ldots,c_{k-1}$) is uniquely defined as its definition is independent of the ordering of the zeros. That $c_j(0)=0$ for all $j$ follows as $\alpha_\ell(0)=0$ for all $\ell$. As the above is the unique way to define a monic polynomial with these zeros (\exerciseref{exercise:monicpolyunique}), the uniqueness part of the theorem follows.

We need to show that the coefficients are holomorphic functions on $V'$, and that $u$ is a holomorphic function on $V$. No matter how you ordered the zeros for each $z'$, the functions $\alpha_\ell$ may not be continuous in general (see \exampleref{sqrt:example}). However, we will prove that the functions $c_\ell$ are holomorphic. The functions $c_\ell$ are (up to sign) the elementary symmetric functions of $\alpha_1,\ldots,\alpha_k$ (see below). A standard theorem in algebra (Newton's identities, see \exerciseref{exercise:powersums}) says that the elementary symmetric functions are polynomials in the so-called power sum functions in the $\alpha_\ell$s:
\[
s_m(z')=\sum_{\ell=1}^k\alpha_\ell(z')^m,\qquad m=1,\ldots,k.
\]
So if the power sums $s_m$ are holomorphic on $V'$, then $c_\ell$ are holomorphic on $V'$. A refinement of the argument principle (see \thmref{thm:onevarargprinc}) says: If $h$ and $g$ are holomorphic functions on a disc $D$, continuous on $\widebar D$, such that $g$ has no zeros on $\partial D$, and $\alpha_1,\ldots,\alpha_k$ are the zeros of $g$ in $D$, then
\[
\frac1{2\pi i}\int_{\partial D}h(\zeta)\frac{g'(\zeta)}{g(\zeta)}\,d\zeta=\sum_{\ell=1}^kh(\alpha_\ell).
\]
With $h(\zeta)=\zeta^m$ and $g(\zeta)=f(z',\zeta)$, the theorem says
\[
s_m(z')=\sum_{\ell=1}^k\alpha_\ell(z')^m=\frac1{2\pi i}\int_{\partial D}\zeta^m\frac{\frac{\partial f}{\partial z_n}(z',\zeta)}{f(z',\zeta)}\,d\zeta.
\]
The function $s_m$ is clearly continuous, and if we differentiate under the integral with $\frac\partial{\partial\bar z_1},\ldots,\frac\partial{\partial\bar z_{n-1}}$ we find that $s_m$ is holomorphic. Thus $c_0,\ldots,c_{k-1}$ are holomorphic, and so $P$ is a Weierstrass polynomial.

Finally, we wish to show that $P$ divides $f$ as claimed, that is, that $u$ is holomorphic. For each fixed $z'$, one-variable theory says that $z_n\mapsto\frac{f(z',z_n)}{P(z',z_n)}$ has only removable singularities, and in fact, it has no zeros as we defined $P$ to exactly cancel them all out. The Cauchy formula on $\nicefrac fP$ then says that the function
\[
u(z',z_n)=\frac1{2\pi i}\int_{\partial D}\frac{f(z',\zeta)}{P(z',\zeta)(\zeta-z_n)}\,d\zeta
\]
is equal to $\frac{f(z',z_n)}{P(z',z_n)}$. The function $u$ is clearly continuous and holomorphic in $z_n$ for each fixed $z'$. Differentiating under the integral shows it is also holomorphic in $z'$.
\end{proof}
